% !PRINTED 779
\weberpage{779}
% running head "Primzahlen in Linearformen." / "779" at the top of this page is not
% transcribed as body text.

Nun ist
\[
\frac{\mu_x}{m} - x = \frac{a-m}{m},
\]
und folglich können wir auf die Grössen $\mu_x$ den Satz \S.\,191, 6.
anwenden, in dem wir $n = x$, $\gamma = 1:m$, $c_n = (a-m):m$ zu setzen
haben. Es ist also
\weq{7}{A(s) = \overset{x}{\underset{0,\,\infty}{\Sigma}} \frac{1}{(mx+a)^s}
= \overset{x}{\underset{1,\,\infty}{\Sigma}} \frac{1}{(mx+a-m)^s},}
so lange $s > 1$ ist, eine stetige Function von $s$, und es ist
\weq{8}{A(s) = \frac{1}{m(s-1)} + C(s),}
worin $C(s)$ für $s=1$ endlich bleibt\footnote{Durch Benutzung bekannter
Sätze über $\Gamma$-Functionen findet man nach \S.\,191, (20):
\[
C(1) = \frac{1}{m} \frac{\Gamma'\left(\frac{a}{m}\right)}{\Gamma\left(\frac{a}{m}\right)}.
\]}.

Wir lassen nun $A$ die sämmtlichen Classen (2) durchlaufen, bezeichnen mit
$\chi$ irgend einen der Charaktere der Gruppe $\mathfrak{N}$ und setzen
\weq{9}{Q(s) = \overset{A}{\Sigma} \chi(A)\, A(s),}
oder, was dasselbe ist,
\weq{10}{Q(s) = \overset{n}{\Sigma} \frac{\chi(n)}{n^s},}
worin sich die Summe auf alle Zahlen $n$ erstreckt. Die Anzahl dieser Summen
$Q$ ist so gross, als die Anzahl der Charaktere, d.\,h. $= \varphi(m)$. Wir
bezeichnen sie, wenn eine Unterscheidung nöthig ist, mit
\[
Q_1, Q_2, \ldots, Q_\mu.
\]

Eine von ihnen, $Q_1$, entspricht dem Hauptcharakter und hat den Ausdruck
\[
Q_1(s) = \overset{n}{\Sigma} \frac{1}{n^s}.
\]

Nun ist [\S.\,11, (21)]:
\[
\overset{A}{\Sigma} \chi(A) = \mu \ \text{oder} = 0,
\]
je nachdem $\chi$ der Hauptcharakter ist oder nicht. Ferner ergiebt sich aus
(8)
%CONT>>
